\section{Solver self-reports}\label{app:solver-claims}
A stored BKS records the checker's verdict on its routes in \id{metadata.standard.evaluation}, as described in Section~\ref{sec:results}. A solver can also state that the incumbent it returns is optimal, for instance when an exact search closes its gap, and the submitter may then record that statement in the optional object \id{metadata.optimality}. Its status \id{solver-reported-optimal} says what kind of statement it is, namely a self-report attributed to the solver that made it. The claim is scoped to the model and profile that the solver names, and the checker neither repeats the search nor verifies the bound behind it. Verification of the routes and the claim of optimality therefore remain two distinct statements, made by two different programs. A record without a claim simply omits the object.

The object has three required fields. Its \id{status} is always \id{solver-reported-optimal}, its \id{target_profile} names the evaluation profile under which the solver claims optimality, and its \id{solver} field identifies the solver together with its version or revision when known. An optional \id{campaign} field may point to the run or to its preserved evidence, and further fields may record the optimum and the dual bound reported by the solver, provided that their values and meaning can be established. Since an optimality claim compares the incumbent with every other feasible solution of the same numerical problem, its target \must{} be the profile under which the stored routes were evaluated. The reference library accordingly refuses to record a new claim whose \id{target_profile} differs from the profile of the evaluated family, for static and TD results alike. Appendix~\ref{app:alternatives} lists the other readings and arithmetic modes that a checker may support.

For example, suppose that a fictitious solver \id{ToyBPC 0.1} solves the companion TD instance of Section~\ref{sec:results} under \id{FleetCostDuration} and reports that the two routes \id{[1]} and \id{[2]}, of exact cost $79/6$, form an optimal solution. The checker verifies these routes as for any submission, and the stored result then carries the following object next to \id{metadata.standard}. The distributed file \path{example-td-result.json} carries no claim, so this excerpt only illustrates the format.
\begin{verbatim}
"optimality": {
  "status": "solver-reported-optimal",
  "target_profile": "td-binary64-exact-v1",
  "solver": "ToyBPC 0.1"
}
\end{verbatim}
From such a record, a reader may conclude that the routes are feasible and cost exactly $79/6$ under the declared benchmark revision, because the checker established this, and that the named solver reported finding no better solution under \id{td-binary64-exact-v1}. The optimality itself rests on the correctness of that solver and of its own arithmetic, which the checker does not audit. A reader who needs a proof therefore treats the claim as a pointer to the evidence of the run. Note that a later submission strictly improving the exact key under the same revision, objective and profile would show the claim to be wrong.

\subsection*{The objective reported with an incumbent}

A solver can report the cost it assigned to a feasible solution under its own search model, whether or not it claims optimality. We record such a report in \id{metadata.solver_report}, where \id{objective} names the reported quantity, \id{value} preserves the report and \id{encoding} tells the reader how to interpret it. An \id{integer} or \id{decimal-text} report keeps the original token as a string, including any printed exponent or trailing zeros, whereas a \id{binary64} report from a solver API uses a finite JSON number with a round-trip spelling. For a known scale, \path{family_units_per_solver_unit} stores the positive rational number of family cost units represented by one solver unit, using the fraction encoding of Section~\ref{sec:results}.

When the report gives a routing cost and its unit is known, \id{metadata.solver_cost} displays that cost in family units. To obtain it, multiply the reported value by the unit exactly and round once to nearest binary64, ties to even. Omit \id{solver_cost} if the routing cost or the conversion is unknown, or if its projection overflows. This display lets a reader set the solver's own view of an incumbent beside the checked value, while the comparison of incumbents always uses the exact cost in \id{metadata.standard}.

For example, an LKH-3 run on a Toronto instance reports the integer \id{"65617528"} for \id{total-arc-cost}. Since its unit is $1/100$, \id{metadata.solver_cost} displays $\RoundBinary(65617528/100)$ as \id{656175.28}. An exact evaluation of the same routes on the family matrix, however, gives $46174228443702207525/70368744177664$, whose binary64 display is \id{656175.252}, and the BKS comparison uses this latter rational. The two values differ because LKH-3 searched on integer-scaled arc costs in units of $1/100$, whereas the checker evaluates the binary64 values of the family matrix exactly.

A solver may also optimize a composite objective. For example, PyVRP can report \id{vehicle-penalized-cost} and separately report the \id{total-arc-cost} of the routes, in which case \path{routing_cost_component} records the latter quantity with its own \id{objective}, \id{value} and \id{encoding}, in the unit of the report, and \id{solver_cost} is computed from this component. Kayros can instead report \id{total-route-duration} directly. An ordinary solver report can accompany any incumbent, whereas a solver that claims optimality identifies the target of that claim separately in \id{metadata.optimality}.
